\chapter{Testing and Results}

Testing was performed to verify that VisuaLearn AI satisfies its functional requirements, responds correctly to learner actions, stores learning resources, and remains usable across devices. The system contains several interacting modules, so testing was carried out at the frontend, backend, database, simulation, visualization, voice, and learning workflow levels.

\section[Testing Strategy]{\textbf{Testing Strategy}}
The testing strategy included the following:
\begin{itemize}[leftmargin=1.2cm]
	\item Functional testing of login, chat, learning resources, simulation, visualization, and settings.
	\item Backend route testing for valid and invalid requests.
	\item Manual testing of complete learner workflows.
	\item Simulation testing using deterministic algorithm examples.
	\item Responsive testing on desktop and mobile screen sizes.
	\item Error handling tests for missing input, unavailable resources, and rendering failures.
	\item Persistence testing for conversations, messages, and learning resources.
\end{itemize}

\section[Functional Test Cases]{\textbf{Functional Test Cases}}
\begin{longtable}{p{1.1cm}p{4cm}p{5.3cm}p{2.3cm}}
\caption{Functional Test Cases}\\
\toprule
\textbf{ID} & \textbf{Scenario} & \textbf{Expected Result} & \textbf{Status}\\
\midrule
TC01 & User login & Valid user is authenticated and redirected to the application & Passed\\
TC02 & New conversation & A new learning session is created and shown in recent sessions & Passed\\
TC03 & Ask question & The system returns an educational response and stores it & Passed\\
TC04 & Open learning workspace & Learning tabs and generated resources are displayed & Passed\\
TC05 & Generate flashcards & Flashcards appear with question and answer structure & Passed\\
TC06 & Generate quiz & Questions, options, and feedback are shown & Passed\\
TC07 & Generate mind map & Concepts are organized as related nodes & Passed\\
TC08 & Detect simulation topic & Suitable procedural topics open simulation workflow & Passed\\
TC09 & Invalid input & A controlled error or fallback response is shown & Passed\\
TC10 & Logout & User session is cleared securely & Passed\\
\bottomrule
\end{longtable}

\section[Learning Workflow Testing]{\textbf{Learning Workflow Testing}}
The learning workflow was tested using topics from computer science, grammar, general science, and engineering. Example queries included merge sort, breadth-first search, parts of speech, machine learning, neural networks, selection sort, and how an engine works. These topics were selected because they require different learning formats.

Algorithmic topics were suitable for simulations and guided steps. Conceptual topics were suitable for explanations, examples, flashcards, and quizzes. Grammar topics were suitable for examples and assessment. This confirmed that the system can organize learning support based on topic type.

\section[Simulation Testing]{\textbf{Simulation Testing}}
Simulation testing focused on state correctness. For sorting algorithms, the array state was checked after each comparison and swap. For graph traversal, the visited node order was inspected. For stack and queue topics, push, pop, enqueue, and dequeue operations were checked. For tree topics, node relationships and traversal order were verified.

The results showed that the simulation module works best for deterministic processes. For topics that do not have clear state transitions, the system avoids forcing simulation and instead provides structured explanations.

\section[Visualization Testing]{\textbf{Visualization Testing}}
Visualization testing verified whether generated visual components load safely and whether fallback messages appear when rendering is not possible. Device capability checks were used to avoid heavy visual rendering on unsupported devices. The isolated widget frame helped prevent invalid visual content from affecting the full application.

\section[Responsive Testing]{\textbf{Responsive Testing}}
The application was tested on desktop and mobile widths. The sidebar, chat input, learning content, flashcard cards, quiz view, simulation view, and settings page were checked for overflow and readability. Mobile improvements focused on keeping the main learning content visible, reducing unnecessary side content, and preventing persona or status labels from overflowing.

\section[Error Handling Testing]{\textbf{Error Handling Testing}}
Error handling was tested by submitting incomplete requests, unavailable resources, unsuitable simulation topics, and rendering failures. The system returned controlled messages instead of exposing raw internal details. Fallback content helped preserve the learning flow even when one feature could not complete.

\section[Observed Results]{\textbf{Observed Results}}
The implementation produced the following results:
\begin{itemize}[leftmargin=1.2cm]
	\item Learners can create and continue learning conversations.
	\item Structured learning resources are generated and displayed in separate views.
	\item Flashcards and quizzes support revision and self-assessment.
	\item Mind maps organize related concepts visually.
	\item Simulations support procedural understanding for algorithmic topics.
	\item Voice interaction supports natural question asking and accessibility.
	\item Persistent storage allows learners to revisit previous sessions.
	\item Adaptive strategy selection supports more personalized response formats.
\end{itemize}

\section[Result Summary]{\textbf{Result Summary}}
Testing confirmed that VisuaLearn AI functions as an integrated adaptive learning workspace. The system performs best when topics can be represented through structured explanation, practice, visualization, or simulation. The platform successfully demonstrates the project objective of moving beyond plain answer generation toward a reusable and interactive learning environment.
